% ─────────────────────────────────────────────────────────────────────────────
%  Guía 2 — Cadenas de Markov (Tiempo Discreto y Continuo)
%  Recopilación, transcripción y resolución de ejercicios · P. Hurtado · 2026
%  Formato de "ficha por ejercicio": enunciado + resolución juntos, con índice
%  navegable (hipervínculos) y diagramas de transición.
% ─────────────────────────────────────────────────────────────────────────────
\documentclass[11pt,a4paper]{article}

\usepackage[T1]{fontenc}
\usepackage[utf8]{inputenc}
\usepackage{lmodern}
\usepackage[a4paper, top=2.0cm, bottom=2.2cm, left=1.8cm, right=1.8cm]{geometry}
\usepackage[spanish,es-noshorthands]{babel}
\usepackage{amsmath,amssymb,mathtools,amsthm}
\usepackage{xcolor}
\usepackage[most]{tcolorbox}
\usepackage{tikz}
\usepackage{pgfplots}
\pgfplotsset{compat=1.18}
\usepackage{microtype}
\usepackage{needspace}
\usepackage{parskip}
\usepackage{fancyhdr}
\usepackage{booktabs}
\usepackage{array}
\usepackage{enumitem}
\usepackage{titlesec}
\usepackage{hyperref}

\usetikzlibrary{arrows.meta, positioning, calc, shapes.geometric, fit, backgrounds, decorations.pathreplacing}

\setlength{\parskip}{4pt}
\setlength{\parindent}{0pt}
\setlength{\headheight}{22pt}
\allowdisplaybreaks
\setcounter{secnumdepth}{0}
\setcounter{tocdepth}{2}

% ─── Paleta (distinta a la guía módulo 03: petróleo + ámbar) ──────────────────
\definecolor{primary}{HTML}{0E5A66}   % teal petróleo
\definecolor{primdark}{HTML}{073B43}
\definecolor{accent}{HTML}{E0913A}    % ámbar
\definecolor{stmtbg}{HTML}{EAF3F4}    % fondo enunciado (teal muy claro)
\definecolor{solbg}{HTML}{FBF4E9}     % fondo resultado (ámbar muy claro)
\definecolor{inkgray}{HTML}{3A3A3A}

\hypersetup{
  colorlinks=true, linkcolor=primary, citecolor=primary, urlcolor=accent,
  pdftitle={Guía 2 — Cadenas de Markov resueltas},
  pdfauthor={P. Hurtado}
}

% ─── Cabeceras ───────────────────────────────────────────────────────────────
\pagestyle{fancy}
\fancyhf{}
\fancyhead[L]{\footnotesize\sffamily\color{primary}Cadenas de Markov \textbullet\ Modelos Estocásticos}
\fancyhead[R]{\footnotesize\sffamily\color{primary}Guía 2 \textbullet\ 2026}
\fancyfoot[C]{\footnotesize\sffamily\color{primary}\thepage}
\renewcommand{\headrulewidth}{0.6pt}
\renewcommand{\headrule}{\hbox to\headwidth{\color{accent}\leaders\hrule height \headrulewidth\hfill}}

% ─── Formato de secciones ────────────────────────────────────────────────────
\titleformat{\section}
  {\Large\bfseries\sffamily\color{primary}}{}{0pt}{}[{\vspace{2pt}\color{accent}\titlerule[2pt]}]
\titlespacing*{\section}{0pt}{16pt}{8pt}
\titleformat{\subsection}
  {\large\bfseries\sffamily\color{primdark}}{}{0pt}{}
\titlespacing*{\subsection}{0pt}{14pt}{4pt}

% ─── Macros matemáticas ──────────────────────────────────────────────────────
\newcommand{\Prob}[1]{\mathbb{P}\!\left(#1\right)}
\newcommand{\E}[1]{\mathbb{E}\!\left[#1\right]}
\newcommand{\R}{\mathbb{R}}

% ─── Caja de enunciado ───────────────────────────────────────────────────────
\newtcolorbox{enun}{
  breakable, enhanced,
  colback=stmtbg, colframe=primary,
  boxrule=0pt, leftrule=4pt, arc=2pt,
  left=10pt, right=10pt, top=18pt, bottom=9pt,
  coltitle=white, fonttitle=\bfseries\sffamily\small,
  colbacktitle=primary, attach boxed title to top left={xshift=8pt, yshift=-9pt},
  boxed title style={colframe=primary, arc=1pt, boxrule=0pt, left=5pt, right=5pt},
  title={ENUNCIADO},
}

% ─── Caja de resultado ───────────────────────────────────────────────────────
\newtcolorbox{resultado}{
  breakable, enhanced,
  colback=solbg, colframe=accent!85!black,
  boxrule=0pt, leftrule=4pt, arc=2pt,
  left=10pt, right=10pt, top=6pt, bottom=6pt,
  fontupper=\normalsize,
}

% ─── Nota al margen / aside ──────────────────────────────────────────────────
\newtcolorbox{nota}{
  breakable, colback=white, colframe=primary!35,
  boxrule=0.6pt, arc=2pt, left=8pt, right=8pt, top=4pt, bottom=4pt,
  fontupper=\footnotesize\itshape\color{inkgray},
}

% ─── Encabezado de paso de la resolución ─────────────────────────────────────
\newcommand{\paso}[1]{%
  \par\addvspace{9pt}\needspace{2.4cm}\noindent
  {\sffamily\bfseries\color{primary}$\blacktriangleright$\ \ #1}\par\nobreak
  \vspace{2pt}\noindent{\color{primary!22}\rule{\linewidth}{0.8pt}}\par\nobreak\vspace{5pt}%
}

% ─── Marca "Resolución" ──────────────────────────────────────────────────────
\newcommand{\resoltag}{%
  \par\addvspace{6pt}\noindent
  \tikz[baseline]{\node[fill=accent,text=white,font=\sffamily\bfseries\footnotesize,
        rounded corners=2pt,inner xsep=6pt,inner ysep=2pt]{RESOLUCIÓN};}\par\vspace{2pt}%
}

% ─── Backlink al índice ──────────────────────────────────────────────────────
\newcommand{\volver}{%
  \par\vspace{6pt}\hfill{\footnotesize\sffamily\hyperlink{indice}{$\uparrow$\ volver al índice}}\par
  \vspace{1pt}{\color{primary!28}\rule{\linewidth}{1.3pt}}\par\addvspace{16pt}%
}

% ─── Estilos TikZ comunes ────────────────────────────────────────────────────
\tikzset{
  edo/.style={circle, draw=primary, fill=primary!10, line width=0.7pt,
              minimum size=0.95cm, font=\small\sffamily},
  edoR/.style={circle, draw=accent!85!black, fill=accent!18, line width=0.9pt,
              minimum size=0.95cm, font=\small\sffamily},
  ar/.style={-{Stealth[length=2.2mm]}, line width=0.7pt, draw=primdark},
  lbl/.style={font=\scriptsize\sffamily, inner sep=1.5pt, fill=white, fill opacity=0.85, text opacity=1},
}

% ═════════════════════════════════════════════════════════════════════════════
\begin{document}
% ═════════════════════════════════════════════════════════════════════════════

\begin{tcolorbox}[
    enhanced, colback=primary, colframe=primdark, colupper=white,
    sharp corners, boxrule=1pt, left=14pt, right=14pt, top=12pt, bottom=12pt,
    borderline west={4pt}{0pt}{accent},
  ]
  {\LARGE\bfseries\sffamily Cadenas de Markov — Ejercicios resueltos}\\[5pt]
  {\normalsize\sffamily Guía 2 \,\textbullet\, Modelos Estocásticos}\\[7pt]
  {\small\sffamily\color{white!85!primary}
    Tiempo discreto (CDMD): clasificación de estados, distribución estacionaria,
    absorción y tiempos de visita. \\
    Tiempo continuo (CDMC): generadores, procesos de nacimiento--muerte y
    probabilidades estacionarias.}\\[7pt]
  {\footnotesize\sffamily\color{white!85!primary}\itshape
    Transcripción y resolución de la carpeta \texttt{guia\_2/} \,\textbullet\, P.~Hurtado \,\textbullet\, 2026}
\end{tcolorbox}

\vspace{4pt}

\begin{nota}
\textbf{Cómo leer esta guía.} Cada ficha contiene el \emph{enunciado} transcrito y, justo
debajo, su \emph{resolución} paso a paso con el o los diagramas de transición. Los
ejercicios se etiquetan \textbf{D\textit{n}} (tiempo \underline{d}iscreto) y \textbf{C\textit{n}}
(tiempo \underline{c}ontinuo). El índice de la página siguiente es navegable: haz clic en
cualquier título para saltar a la ficha; cada ficha tiene al final un enlace
\textsf{$\uparrow$ volver al índice}.
\end{nota}

\clearpage
\hypertarget{indice}{}
{\color{primary}\hrule height 2pt}
\vspace{4pt}
{\Large\bfseries\sffamily\color{primary} Índice navegable}
\vspace{2pt}
{\color{accent}\hrule height 2pt}
\vspace{8pt}
\tableofcontents
\clearpage

% ═════════════════════════════════════════════════════════════════════════════
\section{Parte I · Cadenas en Tiempo Discreto (CDMD)}
% ═════════════════════════════════════════════════════════════════════════════

% ───────────────────────────────────────────── D1
\subsection{D1 · Cadena de 3 estados: probabilidad conjunta}

\begin{enun}
Una cadena de Markov $\{X_n,\ n\ge 0\}$ tiene tres estados $\{0,1,2\}$ y matriz de
transición
\[
  P=\begin{pmatrix} 0.3 & 0.2 & 0.5 \\ 0.6 & 0.3 & 0.1 \\ 0.7 & 0.2 & 0.1\end{pmatrix}.
\]
La distribución inicial es $\Prob{X_0=0}=0.2$, $\Prob{X_0=1}=0.5$, $\Prob{X_0=2}=0.3$.
Determine $\Prob{X_0=0,\ X_1=1,\ X_2=2}$.
\end{enun}

\resoltag
\paso{Regla de la cadena + propiedad de Markov}
La probabilidad de una trayectoria se factoriza como el producto de la inicial por
las transiciones sucesivas:
\[
  \Prob{X_0=0,X_1=1,X_2=2}=\Prob{X_0=0}\,p_{01}\,p_{12}
  = 0.2\times 0.2\times 0.1.
\]
\begin{resultado}
$\displaystyle \Prob{X_0=0,X_1=1,X_2=2}=\boxed{0.004}.$
\end{resultado}
\volver

% ───────────────────────────────────────────── D2
\subsection{D2 · Cadena de 8 estados: análisis estructural}

\begin{enun}
Considere la cadena con estados $\{1,\dots,8\}$ y matriz
\[
  P=\begin{pmatrix}
  0.1&0.2&0.1&0.1&0.2&0.1&0.1&0.1\\
  0&0.1&0.2&0.1&0&0.3&0.1&0.2\\
  0.5&0&0&0.2&0.1&0.1&0.1&0\\
  0&0&0.3&0.7&0&0&0&0\\
  0&0&0.6&0.4&0&0&0&0\\
  0&0&0&0&0&0.3&0.4&0.3\\
  0&0&0&0&0&0.2&0.2&0.6\\
  0&0&0&0&0&0.9&0.1&0
  \end{pmatrix}.
\]
\textbf{a)} Identifique las clases, su tipo (transitoria/recurrente) y el período de
cada estado. \quad \textbf{b)} ¿Existe distribución límite? ¿Y estacionaria?
\end{enun}

\resoltag
\paso{a) Clases, tipo y período}
\emph{Alcanzabilidad.} En el bloque $\{1,2,3,4,5\}$: $1\to2$, $2\to3\to1$,
$1\leftrightarrow3$ ($p_{13},p_{31}>0$), $1\to4\to3\to1$ y $1\to5\to3\to1$. Por
transitividad $\{1,2,3,4,5\}$ es una clase. Las filas $6,7,8$ tienen ceros en las
columnas $1$--$5$: desde $\{6,7,8\}$ no se alcanza $\{1,\dots,5\}$, y $6\leftrightarrow7\leftrightarrow8$.
Luego $\{6,7,8\}$ es una clase \textbf{cerrada}.

\emph{Tipo.} Desde el estado $1$ se alcanza $\{6,7,8\}$ ($p_{16}=0.1>0$) sin retorno
posible $\Rightarrow$ $\{1,\dots,5\}$ es \textbf{transitoria}. $\{6,7,8\}$ cerrada y
finita $\Rightarrow$ \textbf{recurrente positiva}.

\emph{Período.} $p_{44}=0.7>0\Rightarrow d=1$ en $\mathcal C_1$; $p_{66}=0.3>0\Rightarrow d=1$
en $\mathcal C_2$. Ambas clases son \textbf{aperiódicas}.

\begin{center}\small
\begin{tabular}{clll}\toprule
Clase & Estados & Tipo & Período\\\midrule
$\mathcal C_1$ & $\{1,2,3,4,5\}$ & Transitoria & $d=1$\\
$\mathcal C_2$ & $\{6,7,8\}$ & Recurrente positiva & $d=1$\\\bottomrule
\end{tabular}
\end{center}

\paso{b) Distribución límite y estacionaria}
Como $\mathcal C_1$ es transitoria, la cadena es absorbida c.s.\ en $\mathcal C_2$, que es
irreducible, recurrente positiva y aperiódica. Por lo tanto existe \emph{una} distribución
límite, soportada en $\mathcal C_2$:
\[
  \lim_{n\to\infty}P^{n}_{ij}=\pi^{(2)}_j\quad (j\in\mathcal C_2),\qquad
  P^{n}_{ij}\to 0\ \ (j\in\mathcal C_1).
\]
\begin{resultado}
La \textbf{única} distribución estacionaria está concentrada en $\mathcal C_2$ y coincide con
la distribución límite $\boldsymbol\pi^{(2)}$ (solución de $\boldsymbol\pi^{(2)}P_{\mathcal C_2}=\boldsymbol\pi^{(2)}$).
\end{resultado}
\volver

% ───────────────────────────────────────────── D3
\subsection{D3 · Cadena de 5 estados dada por su grafo}

\begin{enun}
Considere la CDMD con estados $\{0,1,2,3,4\}$ dada por el grafo:
\begin{center}
\begin{tikzpicture}
  \node[edo](s0) at (0,2){$0$};
  \node[edo](s1) at (2.6,2){$1$};
  \node[edoR](s2) at (5.2,2){$2$};
  \node[edoR](s3) at (0,0){$3$};
  \node[edoR](s4) at (2.6,0){$4$};
  \draw[ar] (s0) edge[loop above,looseness=6] node[lbl]{$0.5$} (s0);
  \draw[ar] (s3) edge[loop below,looseness=6] node[lbl]{$0.7$} (s3);
  \draw[ar] (s0) -- node[lbl]{$0.5$} (s1);
  \draw[ar] (s1) edge[bend left=15] node[lbl]{$0.5$} (s2);
  \draw[ar] (s1) -- node[lbl]{$0.2$} (s0);
  \draw[ar] (s1) -- node[lbl]{$0.3$} (s4);
  \draw[ar] (s3) -- node[lbl]{$0.3$} (s4);
  \draw[ar] (s4) edge[bend left=20] node[lbl]{$1$} (s3);
\end{tikzpicture}
\end{center}
\textbf{a)} Clases. \textbf{b)} Transitorios/recurrentes (demuestre). \textbf{c)} Periodicidad.
\textbf{d)} $\Prob{X_5=3,X_4=4\mid X_3=1,X_2=0}$. \textbf{e)} Visitas esperadas a $1$ si parte en $1$.
\end{enun}

\resoltag
\paso{a) Clases de comunicación}
$\mathcal C_T=\{0,1\}$: $0\to1$ ($p_{01}=0.5$) y $1\to0$ ($p_{10}=0.2$), se comunican.
$\{2\}$: estado \textbf{absorbente} ($p_{22}=1$), no comunica con ningún otro.
$\mathcal C_R=\{3,4\}$: $3\to4$ ($p_{34}=0.3$) y $4\to3$ ($p_{43}=1$), clase cerrada.
Desde $\mathcal C_T$ se alcanza $\{2\}$ (vía $1\to2$) y $\mathcal C_R$ (vía $1\to4$), pero no hay retorno.

\paso{b) Recurrencia y transiencia}
$\{2\}$ absorbente $\Rightarrow$ \textbf{recurrente positivo}. $\mathcal C_R$ cerrada y finita
$\Rightarrow$ \textbf{recurrente positivo}. $\mathcal C_T$ no es cerrada: desde $1$ se escapa a
$\{2\}$ con prob.\ $0.5$ y a $\{3,4\}$ con prob.\ $0.3$, luego $\mathbf{\{0,1\}\ transitorios}$.

\paso{c) Periodicidad}
$p_{00}=0.5>0\Rightarrow d(0)=1$. $\{2\}$ absorbente: $d(2)=1$.
En $\mathcal C_R$: $p_{33}=0.7>0\Rightarrow d(3)=1$; para $4$, retornos en $2$ y $3$ pasos
($4\to3\to4$ y $4\to3\to3\to4$), $\gcd(2,3)=1\Rightarrow d(4)=1$. Todos \textbf{aperiódicos}.

\paso{d) Probabilidad condicionada}
Por Markov, $X_2=0$ es irrelevante dado $X_3=1$:
\[
  \Prob{X_5=3,X_4=4\mid X_3=1,X_2=0}=p_{14}\,p_{43}=0.3\times 1.
\]
\begin{resultado}$\displaystyle \Prob{\cdot}=\boxed{0.3}.$\end{resultado}

\paso{e) Número esperado de visitas al estado $1$ partiendo en $1$}
Los únicos estados transitorios son $\{0,1\}$; state $2$ es absorbente y $\{3,4\}$ son
recurrentes, ambos contribuyen $0$ a futuras visitas a $1$.
Sea $v_i = \E{\text{\# visitas al estado }1 \mid X_0 = i}$ para $i\in\{0,1\}$.
Por análisis de primer paso:
\begin{alignat}{2}
  v_1 &= 1 + p_{10}\,v_0 &&= 1 + 0.2\,v_0,\label{eq:D3v1}\\
  v_0 &= p_{00}\,v_0 + p_{01}\,v_1 &&= 0.5\,v_0 + 0.5\,v_1.\label{eq:D3v0}
\end{alignat}
\emph{Resolución.}
De \eqref{eq:D3v0}: $0.5\,v_0 = 0.5\,v_1 \;\Rightarrow\; v_0 = v_1$.
Sustituyendo en \eqref{eq:D3v1}:
\[
  v_1 = 1 + 0.2\,v_1 \;\Rightarrow\; 0.8\,v_1 = 1.
\]
\begin{resultado}
$\displaystyle v_1 = \boxed{\tfrac{5}{4} = 1.25\ \text{visitas esperadas al estado }1.}$
\end{resultado}
\volver

% ───────────────────────────────────────────── D4
\subsection{D4 · Cadena de 6 estados: análisis completo}

\begin{enun}
Sea $(X_n)$ una CDMD con $I=\{1,\dots,6\}$ y
\[
  P=\begin{pmatrix}
  0&0.6&0.4&0&0&0\\ 0.2&0&0.8&0&0&0\\ 0.3&0&0.6&0.1&0&0\\
  0&0&0&0&0&1\\ 0&0&0&1&0&0\\ 0&0&0&0&1&0\end{pmatrix}.
\]
\textbf{a)} Grafo y clases. \textbf{b)} Recurrencia/periodicidad. \textbf{c)} $\Prob{X_5=3\mid X_3=1,X_0=3}$
y $\Prob{X_{11}=4\mid X_7\ne1,X_6=2}$. \textbf{d)} Visitas esperadas a $3$ partiendo de $3$.
\textbf{e)} ¿Tiene distribución estacionaria?
\end{enun}

\resoltag
\paso{a) Grafo y clases}
\begin{center}
\begin{tikzpicture}[node distance=1.8cm]
  \node[edo](n1){$1$};
  \node[edo,right=of n1](n2){$2$};
  \node[edo,below=1.3cm of $(n1)!0.5!(n2)$](n3){$3$};
  \node[edoR,right=2.6cm of n2](n4){$4$};
  \node[edoR,below=1.3cm of n4](n6){$6$};
  \node[edoR,left=1.8cm of n6](n5){$5$};
  \draw[ar](n1) edge[bend left=12] node[lbl]{$.6$}(n2);
  \draw[ar](n2) edge[bend left=12] node[lbl]{$.2$}(n1);
  \draw[ar](n1) -- node[lbl]{$.4$}(n3);
  \draw[ar](n2) -- node[lbl]{$.8$}(n3);
  \draw[ar](n3) -- node[lbl]{$.3$}(n1);
  \draw[ar](n3) edge[loop left,looseness=6] node[lbl]{$.6$}(n3);
  \draw[ar](n3) -- node[lbl]{$.1$}(n4);
  \draw[ar](n4) -- node[lbl]{$1$}(n6);
  \draw[ar](n6) -- node[lbl]{$1$}(n5);
  \draw[ar](n5) -- node[lbl]{$1$}(n4);
\end{tikzpicture}
\end{center}
Clases: $\mathcal C_T=\{1,2,3\}$ (se comunican entre sí) y $\mathcal C_R=\{4,5,6\}$
(ciclo $4\to6\to5\to4$). No hay paso de $\mathcal C_R$ a $\mathcal C_T$.

\paso{b) Recurrencia y periodicidad}
$\mathcal C_R$ es cerrada $\Rightarrow$ recurrente positiva. Desde $3$ se escapa a $4$
($p_{34}=0.1$) $\Rightarrow$ $\{1,2,3\}$ \textbf{transitorios}. $p_{33}=0.6>0\Rightarrow d(\mathcal C_T)=1$.
El ciclo $\mathcal C_R$ solo retorna en múltiplos de $3$ $\Rightarrow$ \textbf{período $d=3$}.

\paso{c) Probabilidades}
Por Markov, $\Prob{X_5=3\mid X_3=1,X_0=3}=(P^2)_{13}=p_{12}p_{23}+p_{13}p_{33}
=0.6\cdot0.8+0.4\cdot0.6=\boxed{0.72}$.

Dado $X_6=2$, el único estado con $X_7\ne1$ alcanzable es $X_7=3$ (pues $2\to1$ ó $2\to3$).
Entonces $\Prob{X_{11}=4\mid X_7\ne1,X_6=2}=(P^4)_{34}$. Iterando $\mu^{(k)}=\mu^{(k-1)}P$ desde
$3$:
\[
  \mu^{(1)}=(.3,0,.6,.1,0,0),\ \mu^{(2)}=(.18,.18,.48,.06,0,.1),\ \mu^{(3)}=(.18,.108,.504,.048,.1,.06),
\]
\[
  \mu^{(4)}_4=\mu^{(3)}_3 p_{34}+\mu^{(3)}_5 p_{54}=0.504\cdot0.1+0.1\cdot1=\boxed{0.1504}.
\]

\paso{d) Visitas esperadas a 3}
Con $Q$ sobre $\{1,2,3\}$, $N_{33}=[(I-Q)^{-1}]_{33}$. $\det(I-Q)=0.304$ y cofactor
$M_{33}=1-0.12=0.88$:
\begin{resultado}
$\displaystyle N_{33}=\frac{0.88}{0.304}=\boxed{\tfrac{55}{19}\approx 2.89\ \text{visitas}.}$
\end{resultado}

\paso{e) Distribución estacionaria}
$\mathcal C_R$ es la única clase recurrente: existe estacionaria concentrada en $\{4,5,6\}$,
con $\boldsymbol\pi=(\tfrac13,\tfrac13,\tfrac13)$ sobre el ciclo. Como $d=3$, \emph{no} hay
distribución límite puntual.
\volver

% ───────────────────────────────────────────── D5
\subsection{D5 · Fiabilidad de máquinas}

\begin{enun}
Una máquina está al inicio de cada día \emph{buena} (B) o \emph{fallada} (F). De B:
permanece B con prob.\ $0.98$ y falla con $0.02$. De F (reparación): queda B con
$0.97$ y sigue F con $0.03$.
\textbf{a)} Matriz $P$ de $X_n$ y distribución estacionaria. \textbf{b)} Con $3$ máquinas
i.i.d., $Y_n=\#$ funcionando: matriz de $\{Y_n\}$ y justificación de que es CDMD homogénea.
\end{enun}

\resoltag
\paso{a) Una máquina}
\[
  P=\bordermatrix{&B&F\cr B&0.98&0.02\cr F&0.97&0.03}.
\]
Irreducible y aperiódica $\Rightarrow$ estacionaria única. De $0.02\,\pi_B=0.97\,\pi_F$ y
$\pi_B+\pi_F=1$:
\begin{resultado}
$\displaystyle \pi_B=\tfrac{97}{99}\approx0.98,\qquad \pi_F=\tfrac{2}{99}\approx0.02.$
\end{resultado}

\paso{b) Tres máquinas independientes}
Dado $Y_n=k$, hay $k$ buenas y $3-k$ falladas que transicionan \emph{independientemente};
el futuro depende solo de $k$ y las tasas no cambian con $n$ $\Rightarrow$ \textbf{CDMD homogénea}.
La transición es la convolución de dos binomiales:
\[
  p_{kj}=\!\!\sum_{\substack{r+s=j\\0\le r\le k,\ 0\le s\le 3-k}}\!\!
  \binom{k}{r}0.98^{r}0.02^{k-r}\binom{3-k}{s}0.97^{s}0.03^{3-k-s}.
\]
\begin{center}\small
\begin{tabular}{c|cccc}\toprule
$k$ & $p_{k0}$ & $p_{k1}$ & $p_{k2}$ & $p_{k3}$\\\midrule
0 & $2.7\!\times\!10^{-5}$ & $0.00262$ & $0.08468$ & $0.91267$\\
1 & $1.8\!\times\!10^{-5}$ & $0.00205$ & $0.07585$ & $0.92208$\\
2 & $1.2\!\times\!10^{-5}$ & $0.00156$ & $0.06684$ & $0.93159$\\
3 & $8\!\times\!10^{-6}$ & $0.00118$ & $0.05762$ & $0.94119$\\\bottomrule
\end{tabular}
\end{center}
\volver

% ───────────────────────────────────────────── D6
\subsection{D6 · Inventario de notebooks con demanda Poisson}

\begin{enun}
La demanda $D_t\sim\mathrm{Poisson}(1)$ i.i.d.; $X_t$ es el stock al \emph{final} del
período $t$. Si el stock es $0$ al final del período se ordenan $3$ notebooks (lead time
nulo); si no, no se ordena.
\textbf{a)} Relación $X_{t+1}$ vs $X_t,D_{t+1}$. \textbf{b)} ¿Es CDMD homogénea? \textbf{c)}
Demuestre la matriz dada. \textbf{d)} Recurrencia/período. \textbf{e)} Estacionaria.
\textbf{f)} \% de tiempo con quiebre.
\end{enun}

\resoltag
\paso{a) Recurrencia y b) Markov}
\[
  X_{t+1}=\begin{cases}\max(3-D_{t+1},0)&X_t=0\ \text{(repone a 3)},\\
  \max(X_t-D_{t+1},0)&X_t\ge1.\end{cases}
\]
Dado $X_t$, el siguiente estado depende solo de $D_{t+1}$ (independiente del pasado) y las
probabilidades no dependen de $t$: \textbf{CDMD homogénea}.

\paso{c) Matriz de transición}
Con $\alpha=\Prob{D=0}=e^{-1}$, $\beta=\Prob{D=1}=e^{-1}$, $\gamma=\Prob{D=2}=\tfrac{e^{-1}}2$,
$\delta=\Prob{D\ge3}$: desde $0$ y $3$ es $\max(3-D,0)$ ($p_{\cdot0}=\delta,p_{\cdot1}=\gamma,
p_{\cdot2}=\beta,p_{\cdot3}=\alpha$); desde $1$: $p_{10}=1-\alpha,p_{11}=\alpha$; desde $2$:
$p_{20}=1-2\alpha,p_{21}=\beta,p_{22}=\alpha$. Esto reproduce
\[
  P=\begin{pmatrix}0.080&0.184&0.368&0.368\\0.632&0.368&0&0\\0.264&0.368&0.368&0\\0.080&0.184&0.368&0.368\end{pmatrix}.
\]

\paso{d) Recurrencia y período}
Irreducible (desde $0$ se alcanza todo y todo regresa a $0$) y con autolazos ($p_{00},p_{11},p_{22}>0$):
\textbf{aperiódica}, recurrente positiva.

\paso{e) Estacionaria y f) quiebre}
Finita, irreducible y aperiódica $\Rightarrow$ estacionaria \textbf{única} ($=$ límite),
solución de $\boldsymbol\pi P=\boldsymbol\pi$, $\sum\pi_j=1$ (filas $0$ y $3$ idénticas,
$\pi_0+\pi_3$ actúa como variable compuesta). Hay quiebre cuando $X_t=0$:
\begin{resultado}
proporción de tiempo con quiebre de inventario $=\boxed{\pi_0}$.
\end{resultado}
\volver

% ───────────────────────────────────────────── D7
\subsection{D7 · Buffer de comunicaciones}

\begin{enun}
Buffer de capacidad $4$. Llegadas Poisson tasa $\lambda$ (se pierden si está lleno). Dos
canales: al inicio de cada período se transmiten $\min(X_n,2)$ mensajes; los que llegan
durante el período esperan al siguiente.
\textbf{a)} Recurrencia $X_{n+1}$. \textbf{b)} Matriz (en $p_k=e^{-\lambda}\lambda^k/k!$).
\textbf{c)} Mensajes perdidos por unidad de tiempo en el largo plazo. \textbf{d)} ¿Sirve con
arribos i.i.d.\ no Poisson?
\end{enun}

\resoltag
\paso{a) Recurrencia}
Con $A_n\sim\mathrm{Poisson}(\lambda)$ i.i.d.\ y $R=\max(X_n-2,0)$ mensajes residuales:
\[
  \boxed{\,X_{n+1}=\min\!\big(\max(X_n-2,0)+A_n,\ 4\big).}
\]

\paso{b) Matriz de transición}
Sea $P_{\ge m}=\sum_{k\ge m}p_k$. Según el residuo $R$:
\begin{center}\small
\begin{tabular}{ccc}\toprule
$X_n$ & $R$ & transición\\\midrule
$0,1,2$ & $0$ & $p_{kj}=p_j\ (j\le3),\ p_{k4}=P_{\ge4}$\\
$3$ & $1$ & $p_{3j}=p_{j-1}\ (j=1,2,3),\ p_{34}=P_{\ge3}$\\
$4$ & $2$ & $p_{4j}=p_{j-2}\ (j=2,3),\ p_{44}=P_{\ge2}$\\\bottomrule
\end{tabular}
\end{center}

\paso{c) Mensajes perdidos}
En estado $k$ se pierden $\max(A_n-c_k,0)$ con $c_0=c_1=c_2=4,\ c_3=3,\ c_4=2$:
\[
  \bar L=\sum_{k=0}^{4}\pi_k\,L_k,\qquad L_k=\E{\max(A-c_k,0)}=\sum_{m>c_k}(m-c_k)p_m.
\]

\paso{d) Arribos generales}
Sí: la propiedad de Markov se mantiene mientras $\{A_n\}$ sea i.i.d., pues dado $X_n$ la ley
de $X_{n+1}$ no depende del pasado.
\volver

% ───────────────────────────────────────────── D8
\subsection{D8 · Ayudantes entre dos salas (Solemne)}

\begin{enun}
Dos ayudantes en salas $A$/$B$. Cada $10$ min se lanza un dado: $1$ó$2$ $\Rightarrow$ ambos
cambian; $3$ó$4$ $\Rightarrow$ ayudante 1 cambia; $5$ó$6$ $\Rightarrow$ ayudante 2 cambia.
$X_n=\#$ ayudantes en $A$, estados $\{0,1,2\}$.
\textbf{a)} Matriz y grafo. \textbf{b)} $\Prob{X_2=2\mid X_0=1}$. \textbf{c)} Proporción del
tiempo en la misma sala. \textbf{d)} Variante (otra regla en estados $0,2$): clases, tipos,
períodos.
\end{enun}

\resoltag
\paso{a) Matriz y grafo}
Analizando cada estado (probabilidad $\tfrac13$ por par de caras):
\[
  P=\begin{pmatrix}0&\tfrac23&\tfrac13\\[2pt]\tfrac13&\tfrac13&\tfrac13\\[2pt]\tfrac13&\tfrac23&0\end{pmatrix}.
\]
\begin{center}
\begin{tikzpicture}[node distance=2.4cm]
  \node[edo](a0){$0$};\node[edo,right=of a0](a1){$1$};\node[edo,right=of a1](a2){$2$};
  \draw[ar](a0)edge[bend left=18]node[lbl]{$\tfrac23$}(a1);
  \draw[ar](a1)edge[bend left=18]node[lbl]{$\tfrac13$}(a0);
  \draw[ar](a1)edge[bend left=18]node[lbl]{$\tfrac13$}(a2);
  \draw[ar](a2)edge[bend left=18]node[lbl]{$\tfrac23$}(a1);
  \draw[ar](a0)edge[bend left=40]node[lbl]{$\tfrac13$}(a2);
  \draw[ar](a2)edge[bend left=40]node[lbl]{$\tfrac13$}(a0);
  \draw[ar](a1)edge[loop above,looseness=6]node[lbl]{$\tfrac13$}(a1);
\end{tikzpicture}
\end{center}

\paso{b) $\Prob{X_2=2\mid X_0=1}$}
$(P^2)_{12}=\sum_k p_{1k}p_{k2}=\tfrac13\cdot\tfrac13+\tfrac13\cdot\tfrac13+\tfrac13\cdot0=\boxed{\tfrac29}.$

\paso{c) Misma sala en el largo plazo}
Irreducible y aperiódica. Por simetría $\pi_0=\pi_2$; con $\boldsymbol\pi P=\boldsymbol\pi$ y
$2\pi_0+\pi_1=1$ se obtiene $\boldsymbol\pi=(\tfrac14,\tfrac12,\tfrac14)$.
\begin{resultado}
proporción con ambos en la misma sala $=\pi_0+\pi_2=\boxed{\tfrac12}.$
\end{resultado}

\paso{d) Variante}
La nueva matriz es
$P'=\big(\begin{smallmatrix}\tfrac12&0&\tfrac12\\\tfrac13&\tfrac13&\tfrac13\\\tfrac12&0&\tfrac12\end{smallmatrix}\big)$.
Como $p'_{01}=p'_{21}=0$ pero $p'_{10}>0$, el estado $1$ \textbf{escapa} y no regresa:
$1$ \textbf{transitorio} ($d=1$ por autolazo). $\{0,2\}$ es clase \textbf{recurrente} cerrada,
aperiódica ($p'_{00}>0$).
\volver

% ───────────────────────────────────────────── D9
\subsection{D9 · Deportista y zapatillas en dos puertas}

\begin{enun}
$5$ pares de zapatillas; sale y entra por puerta delantera (D) o de servicio (S), cada una
con prob.\ $\tfrac12$, independientes. Deja las zapatillas en la puerta por la que entra; si
al salir no hay par en su puerta, corre descalzo.
\textbf{a)} Modelo (estado, espacio, $P$). \textbf{b)} Proporción de tiempo descalzo.
\end{enun}

\resoltag
\paso{a) Modelo}
$X_n=\#$ pares en la puerta D al inicio del día (en S hay $5-X_n$), estados $\{0,\dots,5\}$.
Para $k\in\{1,\dots,4\}$ las cuatro combinaciones salida/entrada dan
$p_{k,k-1}=\tfrac14,\ p_{kk}=\tfrac12,\ p_{k,k+1}=\tfrac14$; en las fronteras
$p_{00}=p_{55}=\tfrac34$, $p_{01}=p_{54}=\tfrac14$:
\[
  P=\bordermatrix{&0&1&2&3&4&5\cr
  0&\tfrac34&\tfrac14&&&&\cr
  1&\tfrac14&\tfrac12&\tfrac14&&&\cr
  2&&\tfrac14&\tfrac12&\tfrac14&&\cr
  3&&&\tfrac14&\tfrac12&\tfrac14&\cr
  4&&&&\tfrac14&\tfrac12&\tfrac14\cr
  5&&&&&\tfrac14&\tfrac34}.
\]
\begin{center}
\begin{tikzpicture}[node distance=1.55cm]
  \foreach \i in {0,...,5}{\node[edo](z\i) at (\i*1.7,0){$\i$};}
  \foreach \i [evaluate=\i as \j using int(\i+1)] in {0,...,4}{
     \draw[ar](z\i)edge[bend left=22]node[lbl]{$\tfrac14$}(z\j);
     \draw[ar](z\j)edge[bend left=22]node[lbl]{$\tfrac14$}(z\i);}
  \foreach \i in {1,...,4}{\draw[ar](z\i)edge[loop above,looseness=6]node[lbl]{$\tfrac12$}(z\i);}
  \draw[ar](z0)edge[loop above,looseness=6]node[lbl]{$\tfrac34$}(z0);
  \draw[ar](z5)edge[loop above,looseness=6]node[lbl]{$\tfrac34$}(z5);
\end{tikzpicture}
\end{center}

\paso{b) Proporción descalzo}
Por balance detallado $\pi_k=\pi_{k+1}$ para todo $k$, luego $\boldsymbol\pi$ es \textbf{uniforme},
$\pi_k=\tfrac16$. Sale descalzo si $X_n=0$ y sale por D ($\tfrac12$) o $X_n=5$ y sale por S
($\tfrac12$):
\begin{resultado}
$\displaystyle \Prob{\text{descalzo}}=\pi_0\tfrac12+\pi_5\tfrac12=\boxed{\tfrac16\approx16.7\%.}$
\end{resultado}
\volver

% ───────────────────────────────────────────── D10
\subsection{D10 · Instalación de equipos en un centro hospitalario}

\begin{enun}
Hay $M$ pacientes; para instalar equipos el centro debe estar vacío y \emph{no} ingresan
pacientes hasta lograrlo. Cada mañana cada paciente, indep., sale rehabilitado con prob.\ $p$
y sigue internado con $1-p$.
\textbf{a)} Modele como CDMD, grafo, clases, clasificación. \textbf{b)} Partiendo de $M$:
prob.\ de que algún día haya $M-1$; prob.\ de poder instalar. \textbf{c)} Ya instalados,
llegan $k$ pacientes con prob.\ $q_k\ (k=0,\dots,3)$; sólo $3$ camas; cada cliente paga $C$.
Modifique el modelo, clasifique, plantee estacionaria y dé la ganancia diaria esperada.
\end{enun}

\resoltag
\paso{a) Modelo (sin ingresos)}
$X_n=\#$ pacientes en la semana $n$, estados $\{0,1,\dots,M\}$. Sin llegadas, el número que
\emph{permanece} desde $k$ es $\mathrm{Bin}(k,1-p)$:
\[
  p_{kr}=\binom{k}{r}(1-p)^{r}p^{\,k-r},\qquad 0\le r\le k.
\]
Como sólo se puede bajar ($r\le k$) y hay autolazo ($p_{kk}=(1-p)^k>0$), cada estado
$k\ge1$ es su propia clase \textbf{transitoria}; $\{0\}$ es \textbf{absorbente} (recurrente).
Todos aperiódicos.
\begin{center}
\begin{tikzpicture}[node distance=1.5cm]
  \node[edoR](h0){$0$};\node[edo,right=of h0](h1){$1$};
  \node[edo,right=of h1](h2){$2$};\node[right=1cm of h2](hd){$\cdots$};
  \node[edo,right=1cm of hd](hM){$M$};
  \draw[ar](h1)--node[lbl]{}(h0);
  \draw[ar](h2)edge[bend right=20]node[lbl]{}(h0);
  \draw[ar](h2)--(h1);
  \draw[ar](hd)--(h2);\draw[ar](hM)--(hd);
  \draw[ar](hM)edge[bend left=30]node[lbl]{}(h0);
  \draw[ar](h0)edge[loop above,looseness=6]node[lbl]{$1$}(h0);
  \foreach \n in {h1,h2,hM}{\draw[ar](\n)edge[loop above,looseness=6]node[lbl]{}( \n);}
\end{tikzpicture}
\end{center}

\paso{b) Probabilidades partiendo de $M$}
\emph{Llegar a $M-1$:} sólo es posible en el primer salto que abandona $M$ (después se baja
de $M-1$ sin retorno). Condicionando a salir de $M$:
\[
  \Prob{\text{visitar }M-1}=\frac{p_{M,M-1}}{1-p_{MM}}
  =\frac{M\,p(1-p)^{M-1}}{1-(1-p)^{M}}.
\]
\emph{Poder instalar:} $0$ es absorbente y alcanzable desde todo estado ($p_{k0}=p^{k}>0$);
todos los demás son transitorios $\Rightarrow$ absorción segura:
\begin{resultado}
$\displaystyle \Prob{\text{algún día se instalan los equipos}}=\boxed{1}.$
\end{resultado}

\paso{c) Con ingresos (3 camas)}
Orden diario: parten algunos (cada uno con prob.\ $p$) y luego llegan pacientes hasta llenar
camas. Con $X_n=i$, $S\sim\mathrm{Bin}(i,1-p)$ los que quedan y $A$ los que llegan
($\Prob{A=k}=q_k$), admitidos $=\min(A,3-S)$:
\[
  X_{n+1}=S+\min(A,3-S)=\min(S+A,3),\qquad \text{estados }\{0,1,2,3\}.
\]
\[
  p_{ij}=\sum_{s=0}^{\min(i,j)}\binom{i}{s}(1-p)^{s}p^{\,i-s}\,q_{j-s}\ (j<3),\quad
  p_{i3}=\sum_{s}\binom{i}{s}(1-p)^{s}p^{\,i-s}\Prob{A\ge3-s}.
\]
La cadena es \textbf{irreducible} (sube con llegadas, baja con altas) y \textbf{aperiódica}
(autolazos): finita $\Rightarrow$ estacionaria \textbf{única} $\boldsymbol\pi$, solución de
$\boldsymbol\pi P=\boldsymbol\pi,\ \sum\pi_i=1$. La ganancia diaria es $C$ por admitido:
\begin{resultado}
$\displaystyle \E{\text{ganancia diaria}}=C\sum_{i=0}^{3}\pi_i\,
\E{\min(A,\,3-S)\mid X=i},\quad S\sim\mathrm{Bin}(i,1-p).$
\end{resultado}
\volver

% ───────────────────────────────────────────── D11
\subsection{D11 · Empresa de radiotaxis (Santiago–Concepción)}

\begin{enun}
$3$ automóviles que viajan entre Santiago y Concepción (un viaje diario por auto). Demanda
diaria en Santiago: $0{:}0.25,\,1{:}0.25,\,2{:}0.2,\,3{:}0.1,\,4{:}0.1,\,5{:}0.1$. En
Concepción: $0{:}0.4,\,1{:}0.3,\,2{:}0.3$.
\textbf{(i)} Modele como CDMD: conjunto de estados y probabilidades de transición.
\end{enun}

\resoltag
\paso{(i) Estados y dinámica}
$X_n=\#$ autos en Santiago al inicio del día ($3-X_n$ en Concepción), estados $\{0,1,2,3\}$.
Con demandas independientes $D_S,D_C$, los que parten de Santiago son $\min(X_n,D_S)$ y los
que llegan desde Concepción $\min(3-X_n,D_C)$:
\[
  X_{n+1}=X_n-\min(X_n,D_S)+\min(3-X_n,D_C).
\]
Depende sólo de $X_n$ y de demandas i.i.d.\ $\Rightarrow$ \textbf{CDMD homogénea}. Calculando
cada fila (p.\,ej.\ estado $2$: $\min(2,D_S)\in\{0{:}.25,1{:}.25,2{:}.5\}$,
$\min(1,D_C)\in\{0{:}.4,1{:}.6\}$):
\[
  P=\bordermatrix{&0&1&2&3\cr
  0&0.4&0.3&0.3&0\cr
  1&0.3&0.325&0.3&0.075\cr
  2&0.2&0.4&0.25&0.15\cr
  3&0.3&0.2&0.25&0.25}.
\]
\begin{nota}
La cadena es irreducible y aperiódica (autolazos positivos) $\Rightarrow$ recurrente positiva,
con una única distribución estacionaria $\boldsymbol\pi$ obtenible de $\boldsymbol\pi P=\boldsymbol\pi$.
\end{nota}
\volver

% ───────────────────────────────────────────── D12
\subsection{D12 · El problema de los paraguas}

\begin{enun}
Una persona reparte $r$ paraguas entre su casa y su oficina, caminando entre ambas dos veces
al día. En cada trayecto llueve con probabilidad $p$ (independiente). Si llueve y hay un
paraguas en el lugar de origen, toma uno y lo lleva al destino. Sea $X_n$ el número de
paraguas en el lugar desde donde \emph{parte} en el trayecto $n$.
\textbf{a)} Cadena/transiciones. \textbf{b)} Estacionaria. \textbf{c)} Fracción de trayectos
en que se moja.
\end{enun}

\resoltag
\paso{a) Cadena y matriz}
Estados $\{0,1,\dots,r\}$. Si hay $i>0$ paraguas en el origen: con lluvia ($p$) lleva uno y en
el destino habrá $j=r-i+1$; sin lluvia ($1-p$), $j=r-i$. Si $i=0$, en el destino están los $r$:
\[
  \Prob{i\to j}=\begin{cases}p & j=r-i+1\\ 1-p & j=r-i\end{cases}\ (i>0),\qquad
  \Prob{0\to j}=\begin{cases}1 & j=r\\0&\sim\end{cases}.
\]
\begin{center}
\begin{tikzpicture}[node distance=1.5cm]
  \node[edo](u0){$0$};\node[edo,right=of u0](u1){$1$};\node[edo,right=of u1](u2){$2$};
  \node[right=1cm of u2](ud){$\cdots$};\node[edo,right=1cm of ud](ur){$r$};
  \draw[ar](u0)edge[bend left=30]node[lbl]{$1$}(ur);
  \draw[ar](u1)edge[bend right=20]node[lbl]{$1-p$}(u2.south west);
  \draw[ar](u2)edge[bend right=12]node[lbl]{$1-p$}(ud);
  \draw[ar](ur)edge[bend left=18]node[lbl]{$1-p$}(u0);
  \draw[ar](ur)edge[bend left=10]node[lbl]{$p$}(u1);
\end{tikzpicture}\\[-2pt]
{\scriptsize (la cadena es irreducible y aperiódica: ergódica)}
\end{center}

\paso{b) Distribución estacionaria}
De las ecuaciones de balance se obtiene $\pi_1=\pi_2=\cdots=\pi_r=\dfrac{\pi_0}{1-p}$. Con
$\pi\,(r+1-p)=1$:
\begin{resultado}
$\displaystyle \pi_0=\frac{1-p}{\,r+1-p\,},\qquad \pi_1=\cdots=\pi_r=\frac{1}{\,r+1-p\,}.$
\end{resultado}

\paso{c) Fracción que se moja}
Se moja sólo si está donde no hay paraguas (fracción $\pi_0$) \emph{y} llueve (prob.\ $p$):
\begin{resultado}
$\displaystyle \text{fracción que se moja}=p\,\pi_0=\boxed{\dfrac{p(1-p)}{r+1-p}.}$
\end{resultado}
\volver

% ───────────────────────────────────────────── D13
\subsection{D13 · Clima y seguro de lluvia}

\begin{enun}
El clima diario se modela con estados Despejado (S), Nublado (C) y Lluvioso (R):
\[
  P=\bordermatrix{&S&C&R\cr S&0.5&0.3&0.2\cr C&0.5&0.2&0.3\cr R&0.4&0.5&0.1}.
\]
Hoy está despejado. Un seguro cuesta \$200 y paga \$2500 si los días $7$ y $8$ (contados desde
hoy) son ambos lluviosos. ¿Conviene contratarlo?
\end{enun}

\resoltag
\paso{Distribución a 6 pasos}
Con $f^{(1)}=(1,0,0)^{\!\top}$ (hoy despejado) y $f^{(k)}=(P^{\top})\,f^{(k-1)}$, se evalúa
$f^{(7)}=(P^{\top})^{6}f^{(1)}$:
\[
  f^{(7)}\approx(0.47899,\ 0.31092,\ 0.21009)^{\!\top}
  \ \Rightarrow\ \Prob{X_7=R}\approx0.21009.
\]

\paso{Probabilidad de lluvia en 7 y 8}
\[
  \Prob{X_7=R,X_8=R}=\Prob{X_7=R}\,p_{RR}=0.21009\times0.1=0.021009.
\]

\paso{Decisión por valor esperado}
\[
  \E{\text{pago}}=2500\times0.021009\approx 52.5\ \$ \ <\ 200\ \$.
\]
\begin{resultado}
El pago esperado (\$52.5) es menor que el costo (\$200): \textbf{no conviene contratar el
seguro}.
\end{resultado}
\volver

% ═════════════════════════════════════════════════════════════════════════════
\section{Parte II · Cadenas en Tiempo Continuo (CDMC)}
% ═════════════════════════════════════════════════════════════════════════════

% ───────────────────────────────────────────── C1
\subsection{C1 · Banco de sangre}

\begin{enun}
A un banco de sangre llegan donantes según un Poisson de tasa $\lambda$ (1 litro c/u). Una
ambulancia llega según un Poisson de tasa $\mu$ y retira \emph{toda} la sangre almacenada.
Modele como CDMC y determine las probabilidades estacionarias.
\end{enun}

\resoltag
\paso{Generador}
Estado $=$ litros almacenados $\in\{0,1,2,\dots\}$. Desde $n$: un donante ($\lambda$) lleva a
$n+1$; la ambulancia ($\mu$) lleva a $0$ (desde $0$ no hace nada).
\begin{center}
\begin{tikzpicture}[node distance=1.6cm]
  \node[edo](b0){$0$};\node[edo,right=of b0](b1){$1$};\node[edo,right=of b1](b2){$2$};
  \node[edo,right=of b2](b3){$3$};\node[right=0.9cm of b3](bd){$\cdots$};
  \draw[ar](b0)edge[bend left=18]node[lbl]{$\lambda$}(b1);
  \draw[ar](b1)edge[bend left=18]node[lbl]{$\lambda$}(b2);
  \draw[ar](b2)edge[bend left=18]node[lbl]{$\lambda$}(b3);
  \draw[ar](b3)edge[bend left=18]node[lbl]{$\lambda$}(bd);
  \draw[ar](b1)edge[bend left=30]node[lbl]{$\mu$}(b0);
  \draw[ar](b2)edge[bend left=42]node[lbl]{$\mu$}(b0);
  \draw[ar](b3)edge[bend left=52]node[lbl]{$\mu$}(b0);
\end{tikzpicture}
\end{center}

\paso{Balance y solución}
Para $n\ge1$: $(\lambda+\mu)\pi_n=\lambda\pi_{n-1}\Rightarrow \pi_n=\rho\,\pi_{n-1}$ con
$\rho=\dfrac{\lambda}{\lambda+\mu}$. Luego $\pi_n=\rho^{\,n}\pi_0$ y, normalizando,
$\pi_0=1-\rho=\dfrac{\mu}{\lambda+\mu}$:
\begin{resultado}
$\displaystyle \pi_n=\Big(\frac{\mu}{\lambda+\mu}\Big)\Big(\frac{\lambda}{\lambda+\mu}\Big)^{\!n},
\quad n=0,1,2,\dots\quad(\text{geométrica}).$
\end{resultado}
\volver

% ───────────────────────────────────────────── C2
\subsection{C2 · Lustrabotas de dos sillas}

\begin{enun}
Local con $2$ sillas: en la silla 1 se limpia y embetuna (tasa $\mu_1$) y en la 2 se saca
brillo (tasa $\mu_2$). Clientes potenciales llegan a tasa $\lambda$.
\textbf{a)} Modele cuando el cliente entra \emph{sólo si ambas sillas están vacías}.
Con un ayudante (uno por silla), un cliente entra si la silla 1 está vacía; al terminar la
silla 1 pasa a la 2 si está libre, o espera en la 1.
\textbf{b)} Modele el nuevo problema. \textbf{c)} \% de clientes que entran. \textbf{d)} Tasa
de entrada. \textbf{e)} Número promedio en el local. \textbf{f)} Tiempo promedio dentro.
\end{enun}

\resoltag
\paso{a) Modelo básico (entra sólo si está vacío)}
A lo más un cliente en el sistema. Estados $\{0,\,1,\,2\}$ (vacío / en silla 1 / en silla 2):
\begin{center}
\begin{tikzpicture}[node distance=2.4cm]
  \node[edo](c0){$0$};\node[edo,right=of c0](c1){$1$};\node[edo,right=of c1](c2){$2$};
  \draw[ar](c0)edge node[lbl]{$\lambda$}(c1);
  \draw[ar](c1)edge node[lbl]{$\mu_1$}(c2);
  \draw[ar](c2)edge[bend left=40]node[lbl]{$\mu_2$}(c0);
\end{tikzpicture}
\end{center}

\paso{b) Modelo con ayudante}
Ahora caben hasta $2$ clientes. Estados según ocupación (silla 1, silla 2):
$s_0=(0,0)$, $s_1=(1,0)$, $s_2=(0,1)$, $s_3=(1,1)$ y $s_4=(b,1)$ (silla 1 con cliente
\emph{terminado y bloqueado} esperando, silla 2 ocupada). Es CDMC porque todos los tiempos
(llegadas y servicios) son exponenciales e independientes (sin memoria).
\begin{center}
\begin{tikzpicture}[node distance=2.0cm]
  \node[edo](t0){$s_0$};
  \node[edo,right=of t0](t1){$s_1$};
  \node[edo,below=1.4cm of t1](t2){$s_2$};
  \node[edo,right=2.4cm of t1](t3){$s_3$};
  \node[edo,below=1.4cm of t3](t4){$s_4$};
  \draw[ar](t0)edge node[lbl]{$\lambda$}(t1);
  \draw[ar](t1)edge node[lbl]{$\mu_1$}(t2);
  \draw[ar](t2)edge[bend left=12] node[lbl]{$\mu_2$}(t0);
  \draw[ar](t2)edge node[lbl]{$\lambda$}(t3);
  \draw[ar](t3)edge node[lbl]{$\mu_1$}(t4);
  \draw[ar](t3)edge node[lbl]{$\mu_2$}(t1);
  \draw[ar](t4)edge node[lbl]{$\mu_2$}(t2);
\end{tikzpicture}
\end{center}
Ecuaciones de balance:
\[
\begin{aligned}
\lambda\pi_0&=\mu_2\pi_2, & \mu_1\pi_1&=\lambda\pi_0+\mu_2\pi_3, &
(\mu_2{+}\lambda)\pi_2&=\mu_1\pi_1+\mu_2\pi_4,\\
(\mu_1{+}\mu_2)\pi_3&=\lambda\pi_2, & \mu_2\pi_4&=\mu_1\pi_3, & \textstyle\sum_i\pi_i&=1.
\end{aligned}
\]
Resolviendo en función de $\pi_2$:
\[
  \pi_0=\frac{\mu_2}{\lambda}\pi_2,\quad
  \pi_1=\frac{\mu_2}{\mu_1}\cdot\frac{\mu_1+\mu_2+\lambda}{\mu_1+\mu_2}\,\pi_2,\quad
  \pi_3=\frac{\lambda}{\mu_1+\mu_2}\,\pi_2,\quad
  \pi_4=\frac{\mu_1\lambda}{\mu_2(\mu_1+\mu_2)}\,\pi_2,
\]
y $\pi_2$ se fija con $\sum_i\pi_i=1$.

\paso{c)–f) Medidas de desempeño}
Un cliente entra si la silla 1 está libre (estados $s_0,s_2$). Por PASTA:
\begin{resultado}
\textbf{(c)} $\Prob{\text{entrar}}=\pi_0+\pi_2$.\qquad
\textbf{(d)} tasa de entrada $=\lambda(\pi_0+\pi_2)$.\\[2pt]
\textbf{(e)} $L=\pi_1+\pi_2+2\pi_3+2\pi_4$.\qquad
\textbf{(f)} $W=\dfrac{L}{\lambda(\pi_0+\pi_2)}$ \ (Little).
\end{resultado}
\volver

% ───────────────────────────────────────────── C3
\subsection{C3 · Pedro, el estudiante sociable}

\begin{enun}
Pedro visita a diario a Martita ($50\%$, estadía exp.\ media $20$ min), Juanita ($35\%$, media
$30$ min) o Pepita ($15\%$, media $60$ min). Tras visitar, va a casa y toma siesta (exp.\ media
$2$ h). Al despertar, con prob.\ $0.3$ juega PlayStation (exp.\ media $45$ min) antes de visitar,
y con $0.7$ va directo. Unidad de tiempo: $1$ día.
\textbf{a)} CDMC: estados y tasas. \textbf{b)} Durmiendo: prob.\ de hacer $3$ visitas sin jugar.
\textbf{c)} Si juega: prob.\ de jugar $>1$ h. \textbf{d)} Lleva $15$ min jugando: prob.\ de jugar
$1$ h más y que la siguiente visita sea Martita.
\end{enun}

\resoltag
\paso{a) Estados y tasas instantáneas}
Estados $\{M,J,P,S,G\}$ (visita a cada amiga, siesta, gaming). Tasas de permanencia (en
días$^{-1}$, $1$ día $=1440$ min): $r_M=72,\ r_J=48,\ r_P=24,\ r_S=12,\ r_G=32$. Las tasas
instantáneas $q_{ij}=r_i\times(\text{prob.\ de ruta})$:
\[
\begin{aligned}
&q_{M,S}=72,\quad q_{J,S}=48,\quad q_{P,S}=24,\\
&q_{S,G}=12(0.3)=3.6,\ \ q_{S,M}=12(0.7)(0.5)=4.2,\ \ q_{S,J}=2.94,\ \ q_{S,P}=1.26,\\
&q_{G,M}=32(0.5)=16,\quad q_{G,J}=32(0.35)=11.2,\quad q_{G,P}=32(0.15)=4.8.
\end{aligned}
\]
\begin{center}
\begin{tikzpicture}[node distance=1.7cm]
  \node[edo](S){$S$};
  \node[edo,above left=1.2cm and 2.2cm of S](M){$M$};
  \node[edo,left=2.6cm of S](J){$J$};
  \node[edo,below left=1.2cm and 2.2cm of S](P){$P$};
  \node[edo,right=2.4cm of S](G){$G$};
  \draw[ar](M)edge node[lbl]{$72$}(S);
  \draw[ar](J)edge node[lbl]{$48$}(S);
  \draw[ar](P)edge node[lbl]{$24$}(S);
  \draw[ar](S)edge node[lbl]{$3.6$}(G);
  \draw[ar](S)edge[bend left=10] node[lbl]{$4.2$}(M);
  \draw[ar](G)edge[bend left=20] node[lbl]{$16$}(M);
  \draw[ar](G)edge[bend right=20] node[lbl]{$4.8$}(P);
\end{tikzpicture}
\end{center}

\paso{b) Tres visitas sin jugar}
Cada vez que despierta elige ``directo'' con prob.\ $0.7$, independiente:
\[
  \Prob{3\text{ visitas sin jugar}}=0.7^{3}=\boxed{0.343}.
\]

\paso{c) Jugar más de 1 hora}
$T_G\sim\exp(\text{media }45)$: $\Prob{T_G>60}=e^{-60/45}=e^{-4/3}\approx\boxed{0.264}.$

\paso{d) Memoria nula + ruta}
Por falta de memoria, $\Prob{\text{jugar }60\text{ min más}}=e^{-60/45}\approx0.264$, indep.\ de
la siguiente visita ($\Prob{M}=0.5$):
\begin{resultado}
$\displaystyle \Prob{\cdot}=0.5\times e^{-4/3}\approx\boxed{0.132}.$
\end{resultado}
\volver

% ───────────────────────────────────────────── C4
\subsection{C4 · Armijo, el ejecutivo, y su novia}

\begin{enun}
Armijo recibe e-mails según Poisson tasa $\lambda$; su buzón guarda a lo más $3$ (el resto
rebota) y responde uno a uno en tiempo exp.\ media $1/\mu$. Su novia llama con tiempos exp.\
media $1/\delta$: si la llamada ocurre con $<3$ consultas pendientes, acude de inmediato (exp.\
media $1/\gamma$); si el buzón está lleno, acude apenas termine el mensaje actual. Si la novia
llama \emph{2 veces sin respuesta}, entra en shock y va a retarlo (exp.\ media $1/\beta$),
borrando todos los e-mails. Sin mails pendientes en la oficina, Armijo hace su declaración de
impuestos.
\textbf{a)} Modele la ocupación de Armijo como CDMC. \textbf{b)} Justifique la existencia de
estacionarias y escriba las ecuaciones. \textbf{c)} N.º promedio de ``escenas de celos'' por
hora. \textbf{d)} Fracción del tiempo dedicada a impuestos.
\end{enun}

\resoltag
\begin{nota}
Por la complejidad del enunciado (marcado $(\ast)$) se adopta un espacio de estados que
captura la actividad de Armijo; se indica explícitamente cada supuesto de modelación.
\end{nota}

\paso{a) Estados y transiciones}
\textbf{Estados:} $T$ (en oficina, $0$ mails $\Rightarrow$ declara impuestos); $E_n$ (en oficina
respondiendo, $n=1,2,3$ mails pendientes, sin llamada pendiente); $F$ (buzón lleno con una
llamada \emph{pendiente}: irá donde la novia al terminar el mail actual); $N$ (con la novia,
encuentro normal); $K$ (escena de celos / reto). Supuesto: al volver de $N$ o $K$ el buzón
queda vacío (en $K$ la novia borra los mails).
\[
\begin{array}{ll}
T:\ \xrightarrow{\lambda}E_1,\ \xrightarrow{\delta}N. &
E_1:\ \xrightarrow{\lambda}E_2,\ \xrightarrow{\mu}T,\ \xrightarrow{\delta}N.\\[2pt]
E_2:\ \xrightarrow{\lambda}E_3,\ \xrightarrow{\mu}E_1,\ \xrightarrow{\delta}N. &
E_3:\ \xrightarrow{\mu}E_2,\ \xrightarrow{\delta}F.\\[2pt]
F:\ \xrightarrow{\mu}N,\ \xrightarrow{\delta}K. &
N:\ \xrightarrow{\gamma}T,\qquad K:\ \xrightarrow{\beta}T.
\end{array}
\]
\begin{center}
\begin{tikzpicture}[node distance=1.7cm]
  \node[edo](T){$T$};
  \node[edo,right=of T](E1){$E_1$};
  \node[edo,right=of E1](E2){$E_2$};
  \node[edo,right=of E2](E3){$E_3$};
  \node[edoR,right=of E3](F){$F$};
  \node[edoR,below=1.3cm of E1](N){$N$};
  \node[edoR,below=1.3cm of F](K){$K$};
  \draw[ar](T)edge[bend left=12]node[lbl]{$\lambda$}(E1);
  \draw[ar](E1)edge[bend left=12]node[lbl]{$\mu$}(T);
  \draw[ar](E1)edge[bend left=12]node[lbl]{$\lambda$}(E2);
  \draw[ar](E2)edge[bend left=12]node[lbl]{$\mu$}(E1);
  \draw[ar](E2)edge[bend left=12]node[lbl]{$\lambda$}(E3);
  \draw[ar](E3)edge[bend left=12]node[lbl]{$\mu$}(E2);
  \draw[ar](E3)edge node[lbl]{$\delta$}(F);
  \draw[ar](T)edge node[lbl]{$\delta$}(N);
  \draw[ar](E1)edge node[lbl]{$\delta$}(N);
  \draw[ar](E2)edge node[lbl]{$\delta$}(N);
  \draw[ar](F)edge node[lbl]{$\mu$}(N);
  \draw[ar](F)edge node[lbl]{$\delta$}(K);
  \draw[ar](N)edge[bend left=10]node[lbl]{$\gamma$}(T);
  \draw[ar](K)edge[bend right=20]node[lbl]{$\beta$}(T);
\end{tikzpicture}
\end{center}

\paso{b) Existencia y ecuaciones}
La cadena es finita e irreducible (todo estado comunica con $T$ y viceversa) $\Rightarrow$
\textbf{ergódica}: existe estacionaria única $\boldsymbol\pi$. Balance (tasa de salida $=$ tasa de
entrada):
\[
\begin{aligned}
(\lambda+\delta)\pi_T&=\mu\pi_{E_1}+\gamma\pi_N+\beta\pi_K, &
(\lambda+\mu+\delta)\pi_{E_1}&=\lambda\pi_T+\mu\pi_{E_2},\\
(\lambda+\mu+\delta)\pi_{E_2}&=\lambda\pi_{E_1}+\mu\pi_{E_3}, &
(\mu+\delta)\pi_{E_3}&=\lambda\pi_{E_2},\\
(\mu+\delta)\pi_F&=\delta\pi_{E_3}, &
\gamma\pi_N&=\delta(\pi_T+\pi_{E_1}+\pi_{E_2})+\mu\pi_F,\\
\beta\pi_K&=\delta\pi_F, & \textstyle\sum\pi&=1.
\end{aligned}
\]

\paso{c) Escenas de celos por hora}
Una escena ocurre al entrar a $K$, a tasa $\delta$ desde $F$:
\begin{resultado}
$\displaystyle \text{escenas/hora}=\delta\,\pi_F.$
\end{resultado}

\paso{d) Fracción en impuestos}
Armijo declara impuestos sólo en el estado $T$:
\begin{resultado}
$\displaystyle \text{fracción en impuestos}=\pi_T.$
\end{resultado}
\volver

% ───────────────────────────────────────────── C5
\subsection{C5 · Sucursal bancaria: cajas y depósitos}

\begin{enun}
Llegan personas Poisson tasa $\lambda$. Con prob.\ $0.8$ van a \emph{cajas} ($5$ cajas en
paralelo, cola ilimitada, servicio exp.\ tasa $\mu$) y con $0.2$ a \emph{depósitos} (módulo
único, $5$ sillones de espera; si están ocupados, se va; servicio exp.\ tasa $\beta$).
\textbf{a)} Tasas de nacimiento y muerte del sistema de cajas. \textbf{b)} Con $3$ en la cola de
cajas y $2$ en la de depósitos, prob.\ de que el próximo en terminar servicio sea de cajas.
\end{enun}

\resoltag
\paso{a) Cajas como nacimiento–muerte (M/M/5)}
Las llegadas a cajas son un Poisson adelgazado de tasa $0.8\lambda$. Con $n$ clientes en cajas:
\begin{resultado}
$\displaystyle \lambda_n=0.8\,\lambda\ \ (n\ge0);\qquad
\mu_n=\begin{cases}n\,\mu & 1\le n\le5,\\[2pt] 5\,\mu & n\ge5.\end{cases}$
\end{resultado}
\begin{center}
\begin{tikzpicture}[node distance=1.45cm]
  \foreach \i in {0,...,5}{\node[edo](k\i) at (\i*1.7,0){$\i$};}
  \node[right=0.7cm of k5](kd){$\cdots$};
  \foreach \i [evaluate=\i as \j using int(\i+1)] in {0,...,4}{
     \draw[ar](k\i)edge[bend left=22]node[lbl]{$.8\lambda$}(k\j);}
  \draw[ar](k5)edge[bend left=22]node[lbl]{$.8\lambda$}(kd);
  \foreach \i [evaluate=\i as \j using int(\i+1), evaluate=\j as \m using int(\j)] in {0,...,4}{
     \draw[ar](k\j)edge[bend left=22]node[lbl]{$\m\mu$}(k\i);}
  \draw[ar](kd)edge[bend left=22]node[lbl]{$5\mu$}(k5);
\end{tikzpicture}
\end{center}

\paso{b) ¿Quién termina primero?}
Con $3$ en cola de cajas las $5$ cajas están ocupadas (tasa total $5\mu$); con $2$ en la cola de
depósitos el módulo está ocupado (tasa $\beta$). Los $6$ servicios compiten como exponenciales
independientes; el primero en terminar es de cajas con probabilidad
\begin{resultado}
$\displaystyle \Prob{\text{próximo en terminar}=\text{cajas}}=\frac{5\mu}{5\mu+\beta}.$
\end{resultado}
\volver

% ───────────────────────────────────────────── C6
\subsection{C6 · Creatividad de un artista plástico}

\begin{enun}
Estados: Desborde creativo (D, produce $10$ obras/mes), Meditación (M, $4$ obras/mes) y
Negación (N, $0$). De D sale a N con prob.\ $1/3$ (si no, a M). De N pasa con certeza a M. De M
va equiprobable a N o D. Permanencias exp.\ de media $1,2,1$ meses para D, M, N.
\textbf{a)} CDMC: estados y tasas. \textbf{b)} ¿Es nacimiento–muerte? Calcule las estacionarias.
\textbf{c)} \% de tiempo productivo. \textbf{d)} Tasa de producción en el largo plazo.
\end{enun}

\resoltag
\paso{a) Tasas instantáneas}
Tasas de salida (mes$^{-1}$): $r_D=1,\ r_M=\tfrac12,\ r_N=1$. Con las probabilidades de ruta:
\[
  q_{D,N}=\tfrac13,\ q_{D,M}=\tfrac23,\qquad q_{N,M}=1,\qquad q_{M,D}=\tfrac14,\ q_{M,N}=\tfrac14.
\]
\begin{center}
\begin{tikzpicture}[node distance=2.6cm]
  \node[edo](D){$D$};\node[edo,right=of D](M){$M$};\node[edo,right=of M](N){$N$};
  \draw[ar](D)edge[bend left=14]node[lbl]{$\tfrac23$}(M);
  \draw[ar](M)edge[bend left=14]node[lbl]{$\tfrac14$}(D);
  \draw[ar](M)edge[bend left=14]node[lbl]{$\tfrac14$}(N);
  \draw[ar](N)edge[bend left=14]node[lbl]{$1$}(M);
  \draw[ar](D)edge[bend left=34]node[lbl]{$\tfrac13$}(N);
\end{tikzpicture}
\end{center}

\paso{b) ¿Nacimiento–muerte? Estacionarias}
\textbf{No} es de nacimiento–muerte: desde $D$ se salta tanto a $M$ como a $N$ (no hay un orden
lineal con sólo vecinos). Resolviendo $\boldsymbol\pi Q=0$:
\[
  \pi_D=\tfrac14\pi_M,\qquad \pi_N=\tfrac13\pi_D+\tfrac14\pi_M=\tfrac13\pi_M,\qquad
  \tfrac12\pi_M=\tfrac23\pi_D+\pi_N\ (\text{consistente}).
\]
Con $\pi_D+\pi_M+\pi_N=1$ ($\tfrac{19}{12}\pi_M=1$):
\begin{resultado}
$\displaystyle \boldsymbol\pi=(\pi_D,\pi_M,\pi_N)=\Big(\tfrac{3}{19},\ \tfrac{12}{19},\ \tfrac{4}{19}\Big).$
\end{resultado}

\paso{c) Tiempo productivo}
Produce en $D$ y $M$:
$\ \pi_D+\pi_M=\tfrac{15}{19}\approx\boxed{78.9\%}.$

\paso{d) Tasa de producción a largo plazo}
\[
  10\,\pi_D+4\,\pi_M+0\,\pi_N=\frac{10\cdot3+4\cdot12}{19}=\frac{78}{19}\approx\boxed{4.11\ \text{obras/mes}.}
\]
\volver

\vspace{6pt}
{\color{primary}\hrule height 2pt}
\vspace{3pt}
{\footnotesize\sffamily\color{primary}Fin de la guía \textbullet\ Cadenas de Markov \textbullet\ Modelos Estocásticos 2026}

\end{document}
